%===============================================================================================
% UNIVERSAL PREAMBLE — the visual grammar, shared by all six variants without exception.
%===============================================================================================
% Everything that decides how this document LOOKS lives here. _styles.tex no longer restyles; it
% only flips a small number of ARRANGEMENT switches defined at the bottom of this file.
%
% WHY THAT CHANGED. An earlier revision differentiated the six variants with ornament: shaded
% header bars, rules above and below the label, a short one-sixth-width accent rule, a vertical
% accent bar, filled-square bullets, and one variant with no bullet glyph at all. Held beside the
% six reference resumes this family is modelled on, every one of those marks was an invention —
% the references use ONE header treatment (bold label, thin full-width rule immediately beneath),
% ONE bullet (small and round), and one entry shape. A reader could pick the generated document
% out of the set at a glance, by its decoration. That is the failure this file now prevents.
%
% WHAT LIVES HERE:
%   * the document class and package set. No variant may add a package.
%   * page geometry, and the setting that keeps text extraction clean.
%   * THE VERTICAL RHYTHM, as named lengths, measured off the references.
%   * the section header, the bullet, the entry shape, the skills shape, the masthead — all final.
%   * the arrangement switches, which are the ONLY thing a variant may change.
%
% Assembled in Node, not by LaTeX. Each variant file opens with a preamble marker and a style
% marker that loadResumeTemplates() replaces. A real \input would need these files copied next to
% the .tex in whatever temp directory compileLatexToPdf creates, which is a failure mode with no
% upside.
%
% CONTAINS NO CANDIDATE DATA AND NO SAMPLE CONTENT. Every human-readable string below is either a
% LaTeX control sequence or a fixed section label. See resumeTemplates.test.js.
%===============================================================================================

\documentclass[letterpaper,11pt]{article}

\usepackage{latexsym}
\usepackage[empty]{fullpage}
\usepackage{titlesec}
\usepackage{marvosym}
\usepackage[usenames,dvipsnames]{color}
\usepackage{verbatim}
\usepackage{enumitem}
\usepackage[hidelinks]{hyperref}
\usepackage{fancyhdr}
\usepackage[english]{babel}
\usepackage{tabularx}
\usepackage{fontawesome5}
\usepackage{multicol}
\usepackage{graphicx}
\setlength{\multicolsep}{0pt}
\setlength{\columnsep}{18pt}
\input{glyphtounicode}

\pagestyle{fancy}
\fancyhf{}
\fancyfoot{}
\renewcommand{\headrulewidth}{0pt}
\renewcommand{\footrulewidth}{0pt}

% Page geometry. Every variant shares it: a resume whose margins change between generations reads
% as a different DOCUMENT rather than a different arrangement, and the six are meant to be one
% family.
\addtolength{\oddsidemargin}{-0.6in}
\addtolength{\evensidemargin}{-0.5in}
\addtolength{\textwidth}{1.19in}
\addtolength{\topmargin}{-.7in}
\addtolength{\textheight}{1.4in}

\urlstyle{same}
\raggedbottom
\raggedright
\setlength{\tabcolsep}{0in}

% LEADING, matched to the references rather than left at the class default.
%
% Bullet and skills text is set at \small — 10pt, the same body size all six references use — and
% the class gives 10pt text a 12pt baseline. Measured off the references with pdf.js, their body
% leading is 11.5pt: Priyanka 11.5, Rithish 11.5, Sudeer 11.5. 0.96 puts \small at 11.52pt and
% scales every other size by the same factor, so the relationship between heading leading and body
% leading is preserved.
%
% This is NOT shrinking to buy space. The font size is untouched — 10pt body, 11pt entry headings,
% 12pt section labels — and the floor of 10pt stands. It is worth stating what it buys, though:
% across a full page of roughly 75 lines it recovers about four lines, which for a profile sitting
% just past the end of page one is the difference between a clean single page and a second page
% carrying four lines above a blank three-quarter sheet.
\linespread{0.96}\selectfont

% This is what keeps every one of the six machine-readable: it embeds a ToUnicode map, so a PDF
% text extractor — and therefore an ATS — recovers real characters instead of font glyph indices.
% Never remove it, in any variant.
\pdfgentounicode=1

%-----------------------------------------------------------------------------------------------
% THE VERTICAL RHYTHM — DEFINED ONCE, HERE, AND NOWHERE ELSE.
%-----------------------------------------------------------------------------------------------
% EVERY gap in the document is one of the lengths below, and EXACTLY ONE construct owns each gap.
% No command, section block or style block may emit an ad-hoc \vspace. That rule exists because
% the scheme this replaced used hand-tuned negative \vspace values scattered across the entry
% commands and list ends, which composed unpredictably and printed two education entries on top
% of one another.
%
%   \resumeSectionSep   above a section header          owned by \titlespacing (before)
%   \resumeHeadSep      header rule to first content    owned by \titlespacing (after)
%   \resumeEntrySep     between entries in a section    owned by the section list's itemsep
%   \resumeRowSep       between the two rows of one entry
%   \resumeEntryLead    entry heading to its bullets    owned by the bullet list's topsep
%   \resumeBulletSep    between bullets                 owned by the bullet list's itemsep
%
% Because each gap has one owner, the section lists all set topsep=0pt: the space under a header
% belongs to \titlespacing, and a list adding its own topsep on top is exactly the kind of
% double-count this scheme exists to remove.
%
% THE VALUES ARE MEASURED, NOT CHOSEN. Every reference PDF was read with pdf.js and its glyph
% baselines compared, giving the real baseline-to-baseline distance for each relationship. The
% right-hand column is the same measurement taken from this template's own output, so it says what
% the settings below actually produce rather than what they were meant to produce:
%
% ("to" and not a dash between the numbers: resumeTemplates.test.js flags any run of seven or more
% digits joined only by spaces, dots, dashes or brackets as phone-number-shaped, and a table of
% dashed ranges reads exactly like one to a regex.)
%
%   relationship                      reference spread (all six)   this template
%   ------------------------------    --------------------------   ---------------
%   body leading, 10pt text                11.5 to 12.0              11.5 to 12.0
%   content -> next section header         16.0 to 28.5              19.0 to 21.0
%   section header -> first content        12.5 to 16.5              13.5 to 16.0
%   entry heading -> first bullet          11.5 to 15.0              12.0
%   bullet -> bullet                       11.5 to 15.0              12.0 to 12.5
%   last bullet -> next entry heading      12.0 to 16.0              15.0 to 16.0
%   skills line -> skills line             11.5 to 15.5              12.5
%
% Every one of them lands inside the reference spread. The references are TIGHT: Sudeer and
% Priyanka add essentially nothing anywhere except above a section header, and the loosest of the
% six, Charitha, still only reaches 28.5 there. These values sit near the tight end deliberately,
% because the profiles this generator sees are long enough that the previous settings pushed them
% onto a second page that was then three-quarters empty.
%
% ALL SIX VARIANTS SHARE THESE VALUES. Density is not a differentiation axis: the references are
% one family precisely because their rhythm is the same. A document-level fit may scale all of
% them together — see \resumeScaleRhythm — but no variant may set them individually.
\newlength{\resumeSectionSep}
\newlength{\resumeHeadSep}
\newlength{\resumeEntrySep}
\newlength{\resumeRowSep}
\newlength{\resumeEntryLead}
\newlength{\resumeBulletSep}
\newlength{\resumeDateWidth}
\newlength{\resumeColGap}
\newlength{\resumeNameGap}
\newlength{\resumeRuleGap}

\setlength{\resumeSectionSep}{6pt}
\setlength{\resumeHeadSep}{2pt}
\setlength{\resumeEntrySep}{2pt}
\setlength{\resumeRowSep}{0.5pt}
\setlength{\resumeEntryLead}{1.5pt}
\setlength{\resumeBulletSep}{1pt}
\setlength{\resumeNameGap}{2pt}
% The gap between the header label's baseline and the rule under it. Part of the header treatment
% rather than the document rhythm, but named for the same reason as everything else here.
\setlength{\resumeRuleGap}{1pt}
% The fixed right-hand column every date is set in. Fixed, not \hfill-driven, so the dates in
% Experience, Education, Projects and Extracurricular all land on the SAME right margin.
\setlength{\resumeDateWidth}{1.72in}
% Without a gutter a full-width heading ends flush against its own date, which reads as a
% collision even when the box arithmetic is exactly right.
\setlength{\resumeColGap}{0.16in}

% SCALES THE WHOLE RHYTHM AT ONCE, from the document body, after the variant's style block has had
% its say. This is what lets the generator FIT a resume rather than hope: a profile that spills
% four lines past page one is recompiled tighter until it fits, and one that genuinely needs two
% pages is recompiled airier so the second page carries a real share of the content instead of six
% lines and a blank sheet. See compileFittedResumePdf in latexCompilerService.
%
% It scales the rhythm UNIFORMLY, so a fitted document is still the same family — every gap moves
% by the same factor and the relationships between them are preserved.
%
% \titlespacing is re-issued because titlesec has already read the old values into its own
% internal skips; scaling the lengths alone would leave every section header spaced as before.
\newcommand{\resumeScaleRhythm}[1]{%
  \setlength{\resumeSectionSep}{#1\resumeSectionSep}%
  \setlength{\resumeHeadSep}{#1\resumeHeadSep}%
  \setlength{\resumeEntrySep}{#1\resumeEntrySep}%
  \setlength{\resumeEntryLead}{#1\resumeEntryLead}%
  \setlength{\resumeBulletSep}{#1\resumeBulletSep}%
  \titlespacing*{\section}{0pt}{\resumeSectionSep}{\resumeHeadSep}%
}

% A \needspace equivalent. The bundled TeX distribution does not ship the needspace package and
% the no-new-packages rule stands, so this measures the page directly: how much room is left is
% \pagegoal minus \pagetotal, and if that is less than what the header plus its first content line
% needs, start a new page now so the two travel together.
%
% NOT needspace's own \vskip/\penalty-100/\vskip trick. That works by making a break cheap when
% the page is nearly full — but this document sets \raggedbottom, which parks a \vfil at the foot
% of every page and therefore makes leftover space cost almost nothing no matter how much of it
% there is. Under \raggedbottom the -100 reads as a break BONUS that TeX will take at practically
% any section boundary, and pages started ending a quarter of the way up. Comparing the two
% registers is deterministic and immune to that.
%
% \pagegoal is \maxdimen before any material has been contributed to a page, which is the
% start-of-document case; there is nothing to check then, and comparing against it would always
% fire.
\newcommand{\resumeNeedSpace}[1]{%
  \par
  \begingroup
    \setlength{\dimen0}{#1}%
    \ifdim\pagegoal=\maxdimen\else
      \ifdim\dimen0>\dimexpr\pagegoal-\pagetotal\relax
        \newpage
      \fi
    \fi
  \endgroup
}

%-----------------------------------------------------------------------------------------------
% SECTION HEADER — one treatment, all six variants.
%-----------------------------------------------------------------------------------------------
% Bold label at \large (12pt against an 11pt body, inside the references' 11-13pt band), with a
% thin rule spanning the FULL text width immediately beneath it. That is what all six references
% do and it is not negotiable per variant.
%
% Explicitly not offered, and previously present: a shaded background bar behind the label; a rule
% above as well as below; a short partial-width accent rule; a vertical accent bar beside the
% label; colour of any kind. A reviewer could identify the generated resume in a stack purely by
% those marks.
%
% \titlespacing owns the gap above the header and the gap below the rule; nothing else may add to
% either. The rule is drawn in the after-code so it belongs to the header rather than to the
% content that follows.
%
% THE RULE IS NAMED RATHER THAN WRITTEN INLINE, and that is not a style preference. \titleformat's
% trailing after-code is an optional argument, so LaTeX scans it to the first unmatched `]` — and
% writing `[...\titlerule[0.8pt]]` ends the argument at \titlerule's OWN closing bracket. The
% after-code silently becomes `...\titlerule[0.8pt` and the leftover `]` prints as a literal
% bracket in the top-left corner of page one. Every section rule in the document vanished and a
% stray `]` appeared, from a document that compiled without a single warning.
\newcommand{\resumeSectionRule}{\vspace{\resumeRuleGap}\titlerule[0.8pt]}
\titleformat{\section}{\raggedright\large\bfseries}{}{0em}{}[\resumeSectionRule]
\titlespacing*{\section}{0pt}{\resumeSectionSep}{\resumeHeadSep}

%-----------------------------------------------------------------------------------------------
% BULLET — one glyph, all six variants, every bulleted section.
%-----------------------------------------------------------------------------------------------
% A small round bullet, scaled down from \textbullet and nudged onto the x-height so it sits with
% the text rather than above it. Every reference uses exactly this mark, at roughly this size, in
% Skills, Experience, Projects, Certifications and Coursework alike.
%
% Previously in use and now gone: a filled square, a six-pointed star, a hollow square, a centred
% asterisk, and — in one variant — no marker at all, which rendered bullets as bare indented text
% that read as wrapped prose rather than as a list.
\newcommand{\resumebullet}{\raisebox{0.2ex}{\scalebox{0.7}{\textbullet}}}
% Kept under its old name because previously stored LaTeX references it.
\newcommand{\resumedash}{\resumebullet}

% BULLET LIST under an entry. topsep owns the entry-to-bullets gap; itemsep owns bullet-to-bullet.
% leftmargin is 1.4em to put the text edge ~15pt in from the margin, which is where the references
% hang their wrapped bullet lines.
\newcommand{\resumeItemListStart}{%
  \begin{itemize}[leftmargin=1.4em, label=\resumebullet, align=left, labelwidth=0.85em, labelsep=0.55em,
                  topsep=\resumeEntryLead, partopsep=0pt, parsep=0pt, itemsep=\resumeBulletSep]}
\newcommand{\resumeItemListEnd}{\end{itemize}}
% A single short word alone on a bullet's last line is prevented in the CONTENT builder, which
% ties the final two words together with a non-breaking space. Doing it here with \parfillskip
% would force every short bullet to stretch its last line and report an Underfull \hbox, and
% underfull boxes are part of the acceptance gate.
\newcommand{\resumeItem}[1]{\item \small{#1}}

%-----------------------------------------------------------------------------------------------
% ENTRY ROWS — the two-column line that every dated entry is built from.
%-----------------------------------------------------------------------------------------------
% Two \parbox[t]es of FIXED width, not a tabular*. A tabular* with [t] pins the box's baseline to
% its FIRST row, so every later row became depth that the surrounding spacing happily overprinted;
% and a tabular cell does not wrap, so a long degree name ran straight through its own date and
% out past the margin.
%
% Fixed-width \parbox[t]es fix both at once: [t] aligns the two boxes by their FIRST baselines,
% the row's height is the taller of the two and is fully measured, and long text WRAPS inside its
% own column instead of colliding with the date. The date column is \resumeDateWidth wide in every
% section, which is what makes the right margin identical throughout the document.
\newcommand{\resumeEntryRow}[2]{%
  \noindent
  \parbox[t]{\dimexpr\linewidth-\resumeDateWidth-\resumeColGap\relax}{\raggedright #1}%
  \hspace{\resumeColGap}%
  \parbox[t]{\resumeDateWidth}{\raggedleft \mbox{#2}}%
  \par
}

% Experience / project / extracurricular entry — ONE LINE, all six variants.
%
%     Role | Company, Location                                    Date range
%     (bold, left)                                              (bold, right)
%
% Every reference uses exactly this. The two-line shape that was here before — company bold on
% line one with the role italicised beneath — belongs to a different template family, costs a line
% per role, and was the single most obvious tell that the generated resume was not one of the six.
% The pieces are assembled and joined in the content builder, so a role with no recorded location,
% or an entry with a company but no title, collapses cleanly rather than emitting a stray "|".
%
% \nobreak after the row: an entry heading must never be the last thing on a page with its bullets
% at the top of the next one. The penalty lands in the vertical list immediately before the bullet
% list's topsep glue, and glue preceded by a penalty is not a legal breakpoint, so the heading and
% its first bullet always travel together.
\newcommand{\resumeEntry}[2]{%
  \item
  \resumeEntryRow{\textbf{#1}}{\textbf{#2}}%
  \nobreak
}

% Education entry: #1 degree (bold, left), #2 dates (bold, right), #3 institution, #4 GPA or empty.
% \nobreak around the row gap: the institution line must never be split from its degree by a page
% break.
\newcommand{\resumeEntryWithLine}[4]{%
  \item
  \resumeEntryRow{\textbf{#1}}{\textbf{#2}}%
  \nobreak\vskip\resumeRowSep\nobreak
  \resumeEntryRow{\small #3}{\small #4}%
}

% Project heading: #1 the title (already carrying its own tech stack or link markup), #2 dates.
\newcommand{\resumeProjectHeading}[2]{%
  \item
  \resumeEntryRow{#1}{\textbf{#2}}%
  \nobreak
}
\newcommand{\resumeProjectTech}[1]{%
  \nobreak\vskip\resumeRowSep\nobreak
  \noindent\small{\textbf{Technologies:} #1}\par
}

%-----------------------------------------------------------------------------------------------
% Skills — bulleted, one category per line, label and values on the SAME line.
%-----------------------------------------------------------------------------------------------
% A round bullet, the bold category label, a colon, then the values running on from it. Every one
% of the six references does this and nothing else.
%
% Gone with the ornament: the two-column label table that pushed every value to a fixed indent
% column (which left a ragged trench down the middle of the page and wasted a fifth of the text
% width), the two-column multicol arrangement, and the flowing single paragraph.
%
% ONE LINE PER DEFINITION, deliberately: resumeGeneration.test.js reads these definitions
% line-wise to check that the arrangement is what it claims to be.
\newcommand{\resumeSkillsStart}{\begin{itemize}[leftmargin=1.4em, label=\resumebullet, align=left, labelwidth=0.85em, labelsep=0.55em, topsep=0pt, partopsep=0pt, parsep=0pt, itemsep=\resumeBulletSep]}
\newcommand{\resumeSkillEntry}[2]{\item \small{\textbf{#1:} #2}}
\newcommand{\resumeSkillsEnd}{\end{itemize}}

%-----------------------------------------------------------------------------------------------
% MASTHEAD.
%-----------------------------------------------------------------------------------------------
% Name, then the optional positioning headline, then the contact line with its icon set. The
% headline and contact arguments arrive as \resumeHeadline{...} / \resumeContact{...} calls, or as
% NOTHING AT ALL when the profile has no value — that emptiness decision is made in JavaScript,
% where it is reliable, rather than with a LaTeX \ifx test. It is what lets the block lay itself
% out freely: there is never a stray `\\` left behind by an absent field.
%
% The name is \Large (14.4pt), which is where the references set it. It was \Huge (24.9pt), more
% than twice the body size and far outside the family's range.
% NOTHING FOLLOWS THE MASTHEAD. The gap between the masthead and the first section header is
% \resumeSectionSep, owned by \titlespacing, exactly like the gap above every other section header
% in the document — the one-owner-per-gap rule applies to the first section too.
%
% There used to be a \vskip\resumeHeaderBelow here. It was invisible in the unruled variants (a
% following \addvspace absorbs a preceding plain \vskip rather than adding to it), which is
% precisely why it survived: a length that does nothing looks correct in every test that measures
% the result. It is gone, and the length with it.
\newcommand{\resumeHeadline}[1]{\textbf{\small #1}\par\vskip\resumeNameGap}
\newcommand{\resumeContact}[1]{#1}
% \prevdepth is reset so the FIRST section header is spaced exactly like every other one.
%
% TeX places each box in a vertical list using interline glue computed from \prevdepth — the depth
% of the box above. Ordinary body text is about 2pt deep, but the contact line carries fontawesome
% glyphs that descend around 5pt, and \hrule sets \prevdepth to -1000pt (ignore-depth). Both cases
% shift the header off the rhythm: the void above "Summary" measured 9.8pt of whitespace against
% 8.3pt above every other header in the unruled variants and 11.3pt in the ruled ones, even though
% \titlespacing contributed the identical \resumeSectionSep in all three cases. Declaring a depth
% of zero makes the glue calculation ignore what the masthead happens to be made of, so the space
% above the first header is the same space as above every other header.
\newcommand{\resumeHeader}[3]{%
  {\resumeNameAlign
   {\Large\bfseries #1}\par\vskip\resumeNameGap
   #2#3\par}%
  \resumeHeaderRule
}

%-----------------------------------------------------------------------------------------------
% Coursework — three balanced columns by default.
%-----------------------------------------------------------------------------------------------
% Wrapped in commands rather than written into the shared section block so that the columns-versus-
% inline choice is an arrangement switch a variant can flip, like the others below. \multicolsep is
% zeroed in the preamble so the gap above and below the columns is the section rhythm, not
% multicol's own default.
\newcommand{\resumeCourseworkStart}{\begin{multicols}{3}\begin{itemize}[leftmargin=1.4em, label=\resumebullet, align=left, labelwidth=0.85em, labelsep=0.55em, topsep=0pt, partopsep=0pt, parsep=0pt, itemsep=\resumeBulletSep]}
\newcommand{\resumeCourseItem}[1]{\item \small{#1}}
\newcommand{\resumeCourseworkEnd}{\end{itemize}\end{multicols}}

%-----------------------------------------------------------------------------------------------
% THE ARRANGEMENT SWITCHES — the ONLY things _styles.tex may change.
%-----------------------------------------------------------------------------------------------
% Differentiation comes from arrangement, not ornament. A variant may:
%
%   * order its sections differently          — the variant-*.tex file
%   * carry the positioning headline or not   — the variant-*.tex file
%   * change the skills value separator       — VARIANT_STYLE.skillSeparator, in JavaScript
%   * put a project's tech stack inline or on its own line
%                                             — VARIANT_STYLE.projectTechLayout, in JavaScript
%   * centre or left-align the masthead       — \resumeNameAlign, below
%   * close the masthead with a rule or not   — \resumeHeaderRule, below
%   * run coursework in three columns or as one inline list
%                                             — \resumeCoursework*, above
%
% Nothing else. In particular a variant may NOT redefine \titleformat, \resumebullet,
% \resumeEntry, \resumeSkillEntry or any rhythm length, and resumeGeneration.test.js fails the
% build if one does.
\newcommand{\resumeNameAlign}{\centering}
\newcommand{\resumeHeaderRule}{}
